\documentclass[aspectratio=169,10pt]{beamer}
\usetheme[numbering=fraction,progressbar=frametitle]{metropolis}
\usepackage{fontspec}
\setsansfont{Latin Modern Sans}
\usepackage{amsmath,amssymb,mathtools,graphicx,booktabs}
\usepackage{polyglossia}
\setdefaultlanguage{french}
\setbeamertemplate{navigation symbols}{}
\setbeamerfont{frametitle}{size=\large}
\setbeamerfont{normal text}{size=\normalsize}
\setlength{\parskip}{.45em}
\hypersetup{pdfauthor={},pdfsubject={Lecture slides},pdfkeywords={},pdftitle={Suites de matrices et chaînes de Markov}}
\providecommand{\N}{\mathbb{N}}
\providecommand{\Z}{\mathbb{Z}}
\providecommand{\D}{\mathbb{D}}
\providecommand{\Q}{\mathbb{Q}}
\providecommand{\R}{\mathbb{R}}
\providecommand{\C}{\mathbb{C}}
\providecommand{\vect}[1]{\overrightarrow{#1}}
\providecommand{\vi}{\vec{\imath}}
\providecommand{\vj}{\vec{\jmath}}
\providecommand{\vk}{\vec{k}}
\providecommand{\norm}[1]{\left\lVert #1\right\rVert}
\providecommand{\vc}[2]{\begin{pmatrix}#1\\#2\end{pmatrix}}
\providecommand{\deter}[2]{\det\left(#1,#2\right)}
\providecommand{\abs}[1]{\left\lvert #1\right\rvert}
\providecommand{\dd}{\,\mathrm{d}}
\providecommand{\e}{\mathrm{e}}
\providecommand{\Card}{\operatorname{Card}}
\providecommand{\E}{\operatorname{E}}
\providecommand{\V}{\operatorname{V}}
\providecommand{\degre}{^{\circ}}
\providecommand{\un}[1]{\,\mathrm{#1}}
\providecommand{\sci}[2]{#1\times10^{#2}}
\providecommand{\pgcd}{\operatorname{PGCD}}
\providecommand{\ppcm}{\operatorname{PPCM}}
\title{Suites de matrices et chaînes de Markov}
\subtitle{Cours et exemples}
\date{}
\author{}
\begin{document}
\begin{frame}[plain]
\vfill
{\usebeamerfont{title}\inserttitle\par}
\medskip
{\usebeamerfont{subtitle}\insertsubtitle\par}
\medskip
{\usebeamercolor[fg]{progress bar}\rule{\textwidth}{.5pt}}
\vfill
\end{frame}
\begin{frame}{Objectifs}
Étudier une suite de matrices colonnes définie par \(U_{n+1}=AU_n+B\) ; modéliser une évolution aléatoire à deux ou trois états par un graphe probabiliste et une matrice de transition ; calculer la distribution après \(n\) étapes ; déterminer une distribution invariante ; étudier la convergence dans le cas de deux états.
\end{frame}
\begin{frame}{Activité d'introduction}
Dans un modèle simplifié (données fictives), le temps d'une journée est soit beau (\(S\)), soit pluvieux (\(P\)), et ne dépend que du temps de la veille : après un beau jour, il fait beau avec la probabilité \(0{,}8\) ; après un jour de pluie, il fait beau avec la probabilité \(0{,}4\). Le jour \(0\) est beau. On note \(s_n\) et \(p_n\) les probabilités que le jour \(n\) soit beau ou pluvieux.
\end{frame}
\begin{frame}{Activité d'introduction (2)}
\begin{enumerate}
\setcounter{enumi}{0}\item Représenter la situation par un schéma : deux sommets \(S\) et \(P\), une flèche pour chaque passage possible, portant sa probabilité.

\item À l'aide d'un arbre, calculer \(s_1\), \(p_1\), puis \(s_2\), \(p_2\).

\item Justifier que \(s_{n+1}=0{,}8s_n+0{,}4p_n\) et \(p_{n+1}=0{,}2s_n+0{,}6p_n\).

\end{enumerate}
\end{frame}
\begin{frame}{Activité d'introduction (3)}
\begin{enumerate}
\setcounter{enumi}{3}\item On pose \(\pi_n=\begin{pmatrix}s_n&p_n\end{pmatrix}\) et \(T=\begin{pmatrix}0{,}8&0{,}2\\0{,}4&0{,}6\end{pmatrix}\). Vérifier que \(\pi_{n+1}=\pi_nT\). Calculer \(\pi_{10}\) à la calculatrice. Recommencer en partant d'un jour \(0\) pluvieux. Que constate-t-on ?

\item Déterminer la répartition \(\pi=\begin{pmatrix}a&b\end{pmatrix}\), avec \(a+b=1\), qui vérifie \(\pi T=\pi\). La comparer aux observations de la question 4.

\end{enumerate}
\end{frame}
\begin{frame}{Suites de matrices}
Soit \(A\) une matrice carrée d'ordre \(p\), \(B\) une matrice colonne et \((U_n)\) une suite de matrices colonnes telle que \(U_{n+1}=AU_n+B\) pour tout \(n\). Si \(I_p-A\) est inversible, la matrice \(C=(I_p-A)^{-1}B\) vérifie \(C=AC+B\), et pour tout entier naturel \(n\) :
\[U_n=A^n\left(U_0-C\right)+C.\]
\end{frame}
\begin{frame}{Suites de matrices}
\((I_p-A)C=B\) s'écrit \(C=AC+B\). En soustrayant : \(U_{n+1}-C=A(U_n-C)\). La suite \(V_n=U_n-C\) vérifie \(V_{n+1}=AV_n\), donc \(V_n=A^nV_0\) (récurrence), et \(U_n=A^n(U_0-C)+C\).
\end{frame}
\begin{frame}{Suites de matrices}
Si tous les coefficients de \(A^n\) tendent vers \(0\), alors \(U_n\) converge vers \(C\) (coefficient par coefficient). Ce comportement s'observe à la calculatrice ; on l'admet lorsque l'énoncé le précise.
\end{frame}
\begin{frame}{Chaînes de Markov}
On observe un système à des instants \(0,1,2,\dots\) ; à chaque instant, il se trouve dans un des états \(1,\dots,k\) (\(k=2\) ou \(3\)). On note \(X_n\) l'état à l'instant \(n\). On dit que \((X_n)\) est une \emph{chaîne de Markov} (homogène) si la probabilité de passer de l'état \(i\) à l'état \(j\) ne dépend ni de \(n\) ni des états antérieurs : pour tous \(i\), \(j\) et \(n\) tels que \(P(X_n=i)\neq0\),
\[P_{(X_n=i)}(X_{n+1}=j)=t_{ij}.\]
Les nombres \(t_{ij}\) sont les \emph{probabilités de transition}.
\end{frame}
\begin{frame}{Chaînes de Markov}
La chaîne est représentée par un \emph{graphe probabiliste} : un graphe orienté pondéré dont les sommets sont les états, avec un arc de \(i\) vers \(j\), de poids \(t_{ij}\), lorsque \(t_{ij}>0\). La \emph{matrice de transition} est la matrice carrée \(T=(t_{ij})\). Ses coefficients sont positifs et la somme de chaque ligne vaut \(1\).
\end{frame}
\begin{frame}{Chaînes de Markov}
La \emph{distribution} à l'instant \(n\) est la matrice ligne \(\pi_n=\begin{pmatrix}P(X_n=1)&\cdots&P(X_n=k)\end{pmatrix}\) ; \(\pi_0\) est la distribution initiale.
\end{frame}
\begin{frame}{Chaînes de Markov}
Pour tout entier naturel \(n\) : \(\pi_{n+1}=\pi_nT\) et \(\pi_n=\pi_0T^n\).
\end{frame}
\begin{frame}{Chaînes de Markov}
Par la formule des probabilités totales, avec la partition \(\left(X_n=1\right),\dots,\left(X_n=k\right)\) (on omet les événements de probabilité nulle, qui contribuent \(0\)) :
\[P(X_{n+1}=j)=\sum_{i=1}^{k}P(X_n=i)\,t_{ij},\]
qui est le coefficient \(j\) de \(\pi_nT\). La seconde formule s'obtient par récurrence.
\end{frame}
\begin{frame}{Chaînes de Markov}
Si \(P(X_0=i)=1\), alors pour tout état \(j\), \(P(X_n=j)\) est le coefficient de ligne \(i\) et de colonne \(j\) de \(T^n\). Autrement dit, la probabilité de passer de l'état \(i\) à l'état \(j\) en \(n\) transitions est \(\left(T^n\right)_{ij}\).
\end{frame}
\begin{frame}{Chaînes de Markov}
La distribution initiale est la ligne \(\pi_0\) qui a un \(1\) en position \(i\) et des \(0\) ailleurs. Alors \(\pi_n=\pi_0T^n\) est la ligne \(i\) de \(T^n\).
\end{frame}
\begin{frame}{Chaînes de Markov}
Une distribution \(\pi\) (ligne de coefficients positifs de somme \(1\)) est \emph{invariante} si \(\pi T=\pi\). Si \(\pi_0=\pi\), alors \(\pi_n=\pi\) pour tout \(n\).
\end{frame}
\begin{frame}{Chaînes de Markov}
Soit \(T=\begin{pmatrix}1-a&a\\b&1-b\end{pmatrix}\) avec \(0\le a\le1\), \(0\le b\le1\) et \(0<a+b<2\). La chaîne admet une unique distribution invariante \(\pi=\begin{pmatrix}\dfrac b{a+b}&\dfrac a{a+b}\end{pmatrix}\), et quelle que soit la distribution initiale, \(\pi_n\) converge vers \(\pi\).
\end{frame}
\begin{frame}{Chaînes de Markov}
Écrivons \(\pi_n=\begin{pmatrix}x_n&1-x_n\end{pmatrix}\). Alors \(x_{n+1}=(1-a)x_n+b(1-x_n)=(1-a-b)x_n+b\). Cette suite arithmético-géométrique a pour point fixe \(\ell=\dfrac b{a+b}\), et \(x_n-\ell=(1-a-b)^n(x_0-\ell)\). Comme \(-1<1-a-b<1\), \((1-a-b)^n\to0\) et \(x_n\to\ell\). Une distribution invariante vérifie \(x=(1-a-b)x+b\), donc \(x=\ell\) : elle est unique.
\end{frame}
\begin{frame}{Chaînes de Markov}
On résout le système formé par \(\pi T=\pi\) et la condition \(a+b+c=1\) ; l'une des trois équations de \(\pi T=\pi\) est conséquence des deux autres. La convergence éventuelle de \(\pi_n\) s'observe à la calculatrice.
\end{frame}
\begin{frame}{Chaînes de Markov}
La distribution est une ligne et l'on multiplie à droite par \(T\) : \(\pi_{n+1}=\pi_nT\), et non \(T\pi_n\). Ce sont les lignes de \(T\), et non ses colonnes, qui ont pour somme \(1\).
\end{frame}
\end{document}
